% =========================================================================
% Seguimiento de riesgos — 15-11-2026 (control mensual de noviembre, tanda
% S12, F2). Primer seguimiento del método sin semanas (12-11-2026): revisión
% de los diez riesgos vigilados (R1 a R5, R7, R8, R10, R12 y R14); los cinco
% gestionados (R6, R9, R11, R13 y R15) se revisan al cierre de F2. Plan
% vigente: v8 (12-11-2026). Auditado contra el dev-journal (25-10, 12-11 y
% 15-11) y git. Documento independiente de la memoria; creado copiando
% TEMPLATE.tex (ver README.md de esta carpeta).
% Compilar DESDE ESTA CARPETA (2 pases, sin biber):
%   pdflatex --enable-installer -interaction=nonstopmode -file-line-error
%     -output-directory=build 2026-11-15-S12.tex
% =========================================================================
\documentclass[11pt,a4paper]{article}

% =========================================================================
% Preámbulo COMPARTIDO de los documentos independientes de riesgos
% (docs/thesis/plan-riesgos/ y docs/thesis/seguimiento-riesgos/).
%
% Estos documentos NO forman parte de main.tex (la memoria): son entregables
% sueltos en clase article, con el diseño heredado de los .tex de Gestión de
% Proyectos del autor, pero adaptados a la configuración global de la
% memoria (Carlito, márgenes 2,5 cm, parskip 6 pt, títulos de la guía).
% Un solo preámbulo para toda la familia: un ajuste de estilo aquí se
% propaga a todos los seguimientos futuros (uno por sprint hasta S37).
%
% RUTAS: los \input son relativos a la carpeta del documento, no a esta;
% compilar SIEMPRE desde la carpeta del .tex (plan-riesgos/ o
% seguimiento-riesgos/):
%   pdflatex --enable-installer -interaction=nonstopmode -file-line-error
%     -output-directory=build <documento>.tex        (2 pases; sin biber)
% =========================================================================

% --- Codificación (idéntico a packages.tex de la memoria) ---
\usepackage[utf8]{inputenc}
\usepackage[T1]{fontenc}

% --- Idioma: solo español (aquí no hay resumen trilingüe) ---
% es-noshorthands por coherencia con la memoria: acentos en UTF-8, sin
% caracteres activos (el conflicto con biblatex-apa no aplica aquí, pero
% mantener una única configuración de babel evita sorpresas al copiar texto
% entre la memoria y estos documentos).
\usepackage[spanish, es-tabla, es-noshorthands]{babel}

% --- Márgenes: 2,5 cm laterales e inferior, como la memoria ---
% El margen superior se reparte para alojar la cabecera enmarcada del
% diseño de referencia (2 cm de caja + separación) sin comerse el cuerpo:
% el texto arranca a 4 cm del borde y el ancho de texto sigue siendo 16 cm.
\usepackage[a4paper, left=2.5cm, right=2.5cm, bottom=2.5cm,
            top=1.2cm, includehead, headheight=2.2cm, headsep=0.6cm]{geometry}

% --- Tipografía: Carlito, la misma decisión que en style.tex ---
% (clon métrico libre de Calibri; pdflatex no puede usar la real).
% No se replica el arreglo de versalitas de style.tex: aquí no hay
% bibliografía APA y ningún texto usa \textsc.
\usepackage[sfdefault]{carlito}

% --- Espaciado: interlineado sencillo + 6 pt entre párrafos (guía) ---
\usepackage[skip=6pt]{parskip}

% --- Tablas ---
% longtable: el registro y el plan de respuesta superan una página.
% AVISO ESPERADO: "ignored error: Infinite glue shrinkage" en cada salto de
% página de una longtable (bug conocido e inofensivo, latex3/latex2e#1907;
% ver el comentario largo en packages.tex de la memoria).
\usepackage{array}
\usepackage{longtable}
\usepackage{tabularx}
% \RaggedRight (ragged2e) en vez de \raggedright en columnas muy estrechas:
% el \raggedright de LaTeX desactiva de facto la partición silábica y una
% palabra larga («Identificado») desborda una columna de 1,2 cm; la
% variante de ragged2e mantiene el guionado.
\usepackage{ragged2e}
% Las columnas X se componen como m{...}: el contenido de cada celda queda
% centrado verticalmente (aspecto de formulario del diseño de referencia).
\renewcommand{\tabularxcolumn}[1]{m{#1}}
% Tipos de columna del diseño de referencia:
%   L/C = ancho fijo alineado izquierda/centro; Y/Z = X alineada izq./centro.
% En columnas estrechas se prefiere Y (raggedright) a X justificada: la
% justificación en 3-4 cm abre ríos de blancos.
\newcolumntype{L}[1]{>{\raggedright\arraybackslash}m{#1}}
\newcolumntype{C}[1]{>{\centering\arraybackslash}m{#1}}
\newcolumntype{Y}{>{\raggedright\arraybackslash}X}
\newcolumntype{Z}{>{\centering\arraybackslash}X}
% Filas con aire en todas las tablas (el diseño de referencia lo fijaba
% tabla a tabla; aquí una sola vez).
\renewcommand{\arraystretch}{1.25}

% --- Utilidades ---
\usepackage{xcolor}    % \markfield (resaltado) y \todonote (macros.tex)
\usepackage{amssymb}   % \square, \blacksquare, \circlearrowright
\usepackage{xspace}    % lo necesita \bukit (macros.tex)
\usepackage{enumitem}
\setlist[itemize]{leftmargin=*, topsep=2pt, itemsep=2pt}
\setlist[enumerate]{leftmargin=*, topsep=2pt, itemsep=2pt}

% --- Títulos según la guía de estilo, como en style.tex de la memoria ---
% (sin \chapter: la clase article no lo tiene; el resto, idéntico).
\usepackage{titlesec}
\titleformat{\section}
  {\normalfont\bfseries\fontsize{12pt}{14.4pt}\selectfont\raggedright}
  {\thesection.}{0.5em}{\MakeUppercase}
\titleformat{\subsection}
  {\normalfont\bfseries\fontsize{12pt}{14.4pt}\selectfont}
  {\thesubsection.}{0.5em}{}
\titleformat{\subsubsection}
  {\normalfont\itshape\fontsize{12pt}{14.4pt}\selectfont}
  {\thesubsubsection.}{0.5em}{}

% --- Paginación dinámica (regla estricta del autor, 26-07-2026) ---
% PROHIBIDO el salto de página manual también en estos documentos: el
% espacio lo garantizan estos ganchos (o un \Needspace* local ante un bloque
% indivisible), nunca un \newpage. Los .tex de referencia sí usaban
% \newpage; es la única desviación deliberada respecto a ellos.
\usepackage{needspace}
\AddToHook{cmd/section/before}{\Needspace*{5\baselineskip}}
\AddToHook{cmd/subsection/before}{\Needspace*{5\baselineskip}}
\AddToHook{cmd/subsubsection/before}{\Needspace*{4\baselineskip}}
\widowpenalty=10000
\clubpenalty=10000
\displaywidowpenalty=10000
\raggedbottom
% Estos formularios componen mucho texto en columnas de 2-4 cm: se concede
% a TeX un margen de estirado de emergencia antes de desbordar la línea
% (evita Overfull \hbox menores sin tocar el contenido).
\setlength{\emergencystretch}{1em}

% --- Datos globales y macros de la memoria (única fuente: no duplicar) ---
% =========================================================================
% Variables globales del documento.
% Cambiar un dato AQUÍ lo actualiza en todo el documento (portada,
% metadatos del PDF...). No escribir nunca estos datos "a mano" en otros
% archivos: usar siempre estos comandos.
% =========================================================================

% --- Trabajo ---
\newcommand{\thesisTitle}{BUKIT: red social orientada a objetivos}% TODO: título definitivo
\newcommand{\thesisSubtitle}{Memoria del Trabajo de Fin de Grado}%
\newcommand{\thesisDegree}{Grado en Ingeniería Informática de Gestión y Sistemas de Información}%
\newcommand{\thesisCourse}{2026--2027}% curso académico
\newcommand{\thesisDate}{Bilbao, \today}% TODO: <Lugar, día, mes, año> de la entrega/defensa
\newcommand{\thesisYear}{2027}% año de elaboración y defensa

% --- Personas (la plantilla del centro usa "Apellido1 Apellido2, Nombre") ---
\newcommand{\thesisAuthor}{Liviu Deleanu}% nombre natural (para metadatos del PDF)
\newcommand{\thesisAuthorFormal}{Deleanu, Liviu}% TODO: confirmar formato con dos apellidos
\newcommand{\thesisDirector}{Villamañe Gironés, Mikel}%
\newcommand{\thesisCoDirector}{Apellido1 Apellido2, Nombre}% TODO: solo si hay codirector/a

% --- Universidad ---
\newcommand{\thesisUniversity}{Universidad del País Vasco / Euskal Herriko Unibertsitatea}%
\newcommand{\thesisFaculty}{Escuela de Ingeniería de Bilbao}%

% --- Palabras clave del resumen trilingüe ---
\newcommand{\thesisKeywordsEs}{red social, objetivos, monorepo, TypeScript, PostgreSQL}% TODO
\newcommand{\thesisKeywordsEu}{sare soziala, helburuak, monorepo, TypeScript, PostgreSQL}% TODO: revisar euskera
\newcommand{\thesisKeywordsEn}{social network, goals, monorepo, TypeScript, PostgreSQL}% TODO

% =========================================================================
% Comandos propios. Definir aquí cualquier cosa que se repita en el texto:
% si mañana cambia cómo se escribe, se cambia en un único sitio.
% =========================================================================

% Nombre del producto, siempre con el mismo formato.
\newcommand{\bukit}{\textsc{Bukit}\xspace}

% Términos en inglés, en cursiva de forma consistente: \eng{backend}
\newcommand{\eng}[1]{\textit{#1}}

% Notas de trabajo, visibles en rojo en el PDF. Buscar "todonote" y
% eliminarlas todas antes de entregar.
\newcommand{\todonote}[1]{\textcolor{red}{\textbf{[TODO: #1]}}}


% --- Cabecera y pie (diseño de referencia, sin logos) ---
% La caja de cabecera replica la del .tex de referencia (celda ancha de
% contexto + celda estrecha de identidad). La celda derecha llevaba el logo
% del centro: aquí va texto porque los logos UPV/EHU siguen pendientes de
% la fase de diseño (ver style.tex); cuando existan, sustituir por
% \includegraphics. 11 + 4,9 = 15,9 cm < 16 cm de texto: sin desbordes.
\usepackage{fancyhdr}
\newcommand{\riesgosHeader}{%
  \setlength{\tabcolsep}{0pt}%
  \renewcommand{\arraystretch}{1}%
  \begin{tabular}{|p{11cm}|p{4.9cm}|}
    \hline
    \parbox[c][2cm][c]{11cm}{\centering\bfseries
      {\fontsize{16}{19}\selectfont TRABAJO DE FIN DE GRADO\\[4pt]Curso \thesisCourse}%
    } &
    \parbox[c][2cm][c]{4.9cm}{\centering
      {\fontsize{14}{17}\selectfont\bfseries\bukit}\\[3pt]
      {\small\thesisFaculty}%
    }\\
    \hline
  \end{tabular}%
}
\pagestyle{fancy}
\fancyhf{}
\renewcommand{\headrulewidth}{0pt}
\fancyhead[C]{\riesgosHeader}
\fancyfoot[C]{\thepage}
% La primera página usa estilo "plain": se redefine para que también lleve
% la cabecera (mismo truco que el .tex de referencia).
\fancypagestyle{plain}{%
  \fancyhf{}%
  \renewcommand{\headrulewidth}{0pt}%
  \fancyhead[C]{\riesgosHeader}%
  \fancyfoot[C]{\thepage}%
}

% --- Símbolos y campos de formulario del diseño de referencia ---
\newcommand{\fieldline}[1]{\underline{\hspace{#1}}}
\newcommand{\bigline}{\noindent\rule{\textwidth}{0.4pt}}
\newcommand{\checkbox}{\(\square\)}
\newcommand{\checkedbox}{\(\blacksquare\)}
% Símbolo ⟳ = "si se rellena este apartado, nueva versión del plan".
% El original lo resolvía con un SVG externo (paquete svg + Inkscape);
% aquí basta el símbolo de amssymb: cero dependencias nuevas.
\newcommand{\updatesymbol}{\raisebox{-0.2ex}{\(\circlearrowright\)}}
% Resaltado amarillo del identificador de riesgo en los títulos.
\newcommand{\markfield}[1]{\colorbox{yellow!35}{\strut #1}}

% --- Enlaces y metadatos del PDF (al final, como exige hyperref) ---
% El pdftitle lo fija cada documento con \hypersetup tras cargar esto.
\usepackage[hidelinks]{hyperref}
\hypersetup{pdfauthor={\thesisAuthor}}

\hypersetup{pdftitle={Seguimiento de riesgos — 15-11-2026 (control mensual, S12)}}

\begin{document}

\begin{center}
  {\LARGE\bfseries Seguimiento de riesgos — control mensual de noviembre (S12)}\\[6pt]
\end{center}

\textbf{Fecha:} 15-11-2026 · \textbf{Sprint:} S12 (tanda abierta el 12-11-2026, primera del método sin semanas) · \textbf{Responsable:} \thesisAuthor

Regla de revisión (cap.~7 de la memoria, desde el 12-11-2026): \textbf{vigilados} (R1 a R5, R7, R8, R10, R12 y R14) en cada control mensual; \textbf{gestionados} (R6, R9, R11, R13 y R15) al cierre de cada fase; cualquier riesgo, \textbf{de inmediato}, si se observa su indicador clave.

\textbf{El símbolo \updatesymbol\ indica que, si se rellena ese apartado, el cambio se registra en la tabla de seguimiento del cap.~7 de la memoria y entra en la siguiente versión completa del plan de riesgos} (\texttt{../plan-riesgos/}), que se emite al cierre de cada fase, o de inmediato si un riesgo se materializa, cambia su probabilidad o su impacto, o aparece uno nuevo.

Riesgos revisados:
\textbf{los diez vigilados} en el control mensual de noviembre, adelantado al 15-11 por decisión del autor (el método nuevo se adoptó el 12-11 y el primer control estaba previsto en diciembre).
Los gestionados se revisan al cierre de F2.
Próxima revisión: el control mensual de diciembre, en la primera sesión del mes.

\textbf{Plan emitido a partir de este seguimiento:} ninguno.
No cambia ninguna probabilidad, impacto, medida ni contingencia; R12 sigue activo, sin cambio de estado.

Datos del mes: desde el último seguimiento (25-10) no se entregó ninguna historia; la tanda S12 se abrió el 12-11 con la US-19, sin empezar.
Horas registradas: 17~h~50 en octubre (S7 a S9) y 2~h~30 en noviembre (2~h el 12-11 y 30~min el 15-11), dedicadas al cambio de método, la auditoría de coherencia y este control.
Avance frente al calendario de referencia (cap.~13): F2 lleva entregados 10~SP (US-06, US-07 y US-08) y el calendario preveía 18 a esta fecha (también la US-19, la US-09 y la US-10): unos 8~SP, tres semanas de referencia, por detrás.
Equivalencia observada en F2: 10~SP en unas 18~h, 1,8~h/SP con la memoria y el método incluidos; se mantiene la de planificación, 2,0~h/SP (decisión del autor de este control).

\bigline

\section{Medidas de mitigación (prevención)}

Para cada medida de los riesgos que tocan revisión en este seguimiento: ¿se está aplicando?
Si no, justificar y decidir si se mantiene/modifica/elimina.

% Bloques de 3-4 filas (tabularx no parte entre páginas; receta del S2).
\Needspace*{24\baselineskip}
\underline{\textbf{Riesgos vigilados (R1--R4)}}

\begin{tabularx}{\textwidth}{|C{0.9cm}|Y|Y|C{2.2cm}|L{2.4cm}|}
\hline
\multicolumn{1}{|C{0.9cm}|}{\textbf{Id}} &
\multicolumn{1}{>{\centering\arraybackslash}X|}{\textbf{Medida de mitigación}} &
\multicolumn{1}{>{\centering\arraybackslash}X|}{\textbf{¿Aplicada?}} &
\multicolumn{1}{C{2.2cm}|}{\textbf{Justificación (del NO)}} &
\multicolumn{1}{C{2.4cm}|}{\textbf{Decisión}} \\
\hline
R1 & Toda idea nueva entra por el backlog. & \checkedbox\ Sí (sin funciones nuevas en el mes; la tanda S12 tomó la siguiente historia del backlog, la US-19).\newline \checkbox\ No & — & \checkedbox\ Mantener\newline \checkbox\ Modificar\newline \checkbox\ Eliminar \\
\hline
R2 & Medir velocidad y recalibrar en cada control mensual. & \checkedbox\ Sí (equivalencia de F2 medida en este control: 1,8~h/SP; se mantiene 2,0).\newline \checkbox\ No & — & \checkedbox\ Mantener\newline \checkbox\ Modificar\newline \checkbox\ Eliminar \\
\hline
R3 & Diario al cierre de cada sesión; memoria con hueco; auditoría de coherencia a demanda y antes de cada cierre de fase. & \checkedbox\ Sí (los caps.~1, 7, 8, 10, 13 y 14 y el preliminar «Empleo de IA», al día con el cambio de método; auditoría de coherencia lanzada el 12-11 y aplicada el 15-11).\newline \checkbox\ No & — & \checkedbox\ Mantener\newline \checkbox\ Modificar\newline \checkbox\ Eliminar \\
\hline
R4 & Citar al escribir; pendientes con nota visible. & \checkedbox\ Sí (sin fuentes nuevas en el mes).\newline \checkbox\ No & — & \checkedbox\ Mantener\newline \checkbox\ Modificar\newline \checkbox\ Eliminar \\
\hline
\end{tabularx}

\Needspace*{26\baselineskip}
\underline{\textbf{Riesgos vigilados (R5, R7, R8 y R10)}}

\begin{tabularx}{\textwidth}{|C{0.9cm}|Y|Y|C{2.2cm}|L{2.4cm}|}
\hline
\multicolumn{1}{|C{0.9cm}|}{\textbf{Id}} &
\multicolumn{1}{>{\centering\arraybackslash}X|}{\textbf{Medida de mitigación}} &
\multicolumn{1}{>{\centering\arraybackslash}X|}{\textbf{¿Aplicada?}} &
\multicolumn{1}{C{2.2cm}|}{\textbf{Justificación (del NO)}} &
\multicolumn{1}{C{2.4cm}|}{\textbf{Decisión}} \\
\hline
R5 & Revisar todo antes de integrar; nada sin comprobaciones. & \checkbox\ Sí\newline \checkedbox\ No, en parte: la lectura de \texttt{apps/web/src} por el autor, acción del S9, sigue pendiente. & No hubo sesiones de producto en el mes; la acción se mantiene antes de la US-19. & \checkedbox\ Mantener\newline \checkbox\ Modificar\newline \checkbox\ Eliminar \\
\hline
R7 & Minimización desde el diseño (RGPD). & \checkedbox\ Sí (el modelo de datos no cambió; la auditoría de coherencia retiró de la memoria y del diario dos datos privados del autor).\newline \checkbox\ No & — & \checkedbox\ Mantener\newline \checkbox\ Modificar\newline \checkbox\ Eliminar \\
\hline
R8 & Horas por sesión en el diario; el control mensual compara el avance con el calendario. & \checkedbox\ Sí (primera comparación en este control: tres semanas de retraso).\newline \checkbox\ No & — & \checkedbox\ Mantener\newline \checkbox\ Modificar\newline \checkbox\ Eliminar \\
\hline
R10 & Margen deliberado + documentación continua. & \checkedbox\ Sí (el diario permitió retomar el 12-11 tras diecisiete días sin sesiones).\newline \checkbox\ No & — & \checkedbox\ Mantener\newline \checkbox\ Modificar\newline \checkbox\ Eliminar \\
\hline
\end{tabularx}

\Needspace*{20\baselineskip}
\underline{\textbf{Riesgos vigilados (R12 y R14)}}

\begin{tabularx}{\textwidth}{|C{0.9cm}|Y|Y|C{2.2cm}|L{2.4cm}|}
\hline
\multicolumn{1}{|C{0.9cm}|}{\textbf{Id}} &
\multicolumn{1}{>{\centering\arraybackslash}X|}{\textbf{Medida de mitigación}} &
\multicolumn{1}{>{\centering\arraybackslash}X|}{\textbf{¿Aplicada?}} &
\multicolumn{1}{C{2.2cm}|}{\textbf{Justificación (del NO)}} &
\multicolumn{1}{C{2.4cm}|}{\textbf{Decisión}} \\
\hline
R12 & Capítulos a revisión conforme se cierran; versión vigente cuando haya cambios que merezcan revisión, según la disponibilidad de ambos. & \checkedbox\ Sí (sin entregas forzadas; el trabajo no depende de la revisión).\newline \checkbox\ No & — & \checkedbox\ Mantener (sigue activo: ver ficha)\newline \checkbox\ Modificar\newline \checkbox\ Eliminar \\
\hline
R14 & Auditoría en la CI + Dependabot con triaje semanal y alertas + casos de abuso en la Definición de Hecho. & \checkedbox\ Sí (triaje del 12-11: tres propuestas fusionadas, una cerrada y una aplazada a EP04; \texttt{pnpm audit} sin altas).\newline \checkbox\ No & — & \checkedbox\ Mantener\newline \checkbox\ Modificar\newline \checkbox\ Eliminar \\
\hline
\end{tabularx}

\Needspace*{8\baselineskip}
\underline{\textbf{Cambios realizados en las medidas de mitigación}} \updatesymbol

\begin{tabularx}{\textwidth}{|C{1.4cm}|Y|Y|Y|}
\hline
\multicolumn{1}{|C{1.4cm}|}{\textbf{Id}} &
\multicolumn{1}{>{\centering\arraybackslash}X|}{\textbf{Medida original}} &
\multicolumn{1}{>{\centering\arraybackslash}X|}{\textbf{Cambio realizado}} &
\multicolumn{1}{>{\centering\arraybackslash}X|}{\textbf{Justificación}} \\
\hline
— & Ninguno en este seguimiento. & — & — \\
\hline
\end{tabularx}

\Needspace*{8\baselineskip}
\underline{\textbf{Nuevas medidas de mitigación a añadir}} \updatesymbol

\begin{tabularx}{\textwidth}{|C{1.4cm}|Y|Y|}
\hline
\multicolumn{1}{|C{1.4cm}|}{\textbf{Id}} &
\multicolumn{1}{>{\centering\arraybackslash}X|}{\textbf{Nueva medida}} &
\multicolumn{1}{>{\centering\arraybackslash}X|}{\textbf{Justificación}} \\
\hline
— & Ninguna en este seguimiento. & — \\
\hline
\end{tabularx}

\bigline

\section{Riesgos a revisar en este seguimiento}

Los riesgos \textbf{sin novedad y sin observación relevante} se agrupan en la tabla siguiente (indicador en su valor esperado, sin salto, estado «sin novedad»); el que dejó señal en el periodo (R12) conserva su ficha completa a continuación.

\Needspace*{26\baselineskip}
\begin{tabularx}{\textwidth}{|C{0.9cm}|Y|C{2.6cm}|}
\hline
\multicolumn{1}{|C{0.9cm}|}{\textbf{Id}} &
\multicolumn{1}{>{\centering\arraybackslash}X|}{\textbf{Indicador: valor real (esperado 0 en todos)}} &
\multicolumn{1}{C{2.6cm}|}{\textbf{Nuevo estado}} \\
\hline
R1 & 0 — ninguna tanda ganó puntos tras su planificación y nada entró fuera del backlog. & Sin novedad \\
\hline
R2 & Indicador (tres primeros sprints) ya evaluado en S3. F2: 1,8~h/SP observadas frente a 2,0 de planificación; se mantiene 2,0, porque la web y las épicas de mapa y tiempo real son terreno nuevo. & Sin novedad \\
\hline
R3 & 0 — capítulos al día en la sesión del cambio; diario escrito en cada sesión. & Sin novedad \\
\hline
R4 & 0 — sin afirmaciones nuevas de fuentes externas. & Sin novedad \\
\hline
R5 & Observación: la acción del S9 (leer \texttt{apps/web/src} antes de la US-19) sigue pendiente; el indicador no salta porque no se ha integrado nada nuevo. & Sin novedad \\
\hline
R7 & 0 — el modelo de datos no cambió. & Sin novedad \\
\hline
R8 & Observación: tres semanas de retraso frente al calendario de referencia, por debajo del mes del indicador. Si la tanda S12 no avanza antes de mediados de diciembre, el indicador saltará en el próximo control. & Sin novedad \\
\hline
R10 & 0 — el intervalo más largo sin sesiones fue de diecisiete días (26-10 a 11-11), por debajo del mes del indicador. & Sin novedad \\
\hline
R14 & 0 — sin vulnerabilidades altas; una moderada de desarrollo (\texttt{minimatch}, vía ESLint) que no corta la CI. Cuatro propuestas nuevas de Dependabot para el triaje de la semana siguiente. & Sin novedad \\
\hline
\end{tabularx}

% ------------------ R12 ------------------
\Needspace*{30\baselineskip}
\subsection*{Riesgo: \markfield{R12 — Ciclos de revisión con el tutor} (activo desde S9)}

\textbf{Indicadores (clave y de resolución)}

\begin{tabularx}{\textwidth}{|Z|Z|Z|Z|Z|}
\hline
\textbf{Indicador} & \textbf{Valor esperado} & \textbf{Valor real} & \textbf{¿Cumplido?} & \textbf{Acción / observaciones} \\
\hline
Consulta/capítulo >2 semanas sin respuesta, o revisión con cambios estructurales sobre capítulo cerrado. & 0 & Sigue cumplido: la entrega del 05-10 lleva cuarenta y un días sin respuesta del tutor. & \checkedbox\ Sí\quad \checkbox\ No & La contingencia sigue en marcha: el trabajo no depende de la revisión. Si en el control de diciembre continúa sin respuesta, se reevalúa la probabilidad, lo que obligaría a reemitir el plan. \\
\hline
\end{tabularx}

\textbf{Nueva probabilidad (1–3):} \underline{\ 1\ } · \textbf{Nuevo impacto (1–3):} \underline{\ 2\ } · \textbf{Nueva severidad:} \checkbox\ Crítico · \checkbox\ Significativo · \checkedbox\ Moderado\ \updatesymbol\ (sin cambio)

\textbf{¿Cerrar el riesgo?} \checkbox\ Sí · \checkedbox\ No\ \updatesymbol

\textbf{¿Se ha materializado?} \checkbox\ Sí · \checkedbox\ No\ \updatesymbol

\textbf{Nuevo estado} (cap.~7): \checkbox\ Identificado · \checkbox\ Sin novedad · \checkedbox\ Activo · \checkbox\ Materializado · \checkbox\ Resuelto · \checkbox\ Cerrado

\textbf{Plan de contingencia (solo si se ha materializado)}: no aplicable; la demora todavía no retrasa ningún cierre.

\Needspace*{10\baselineskip}
\bigline

\section{Reflexión desde el último seguimiento}

\textbf{¿Ha ocurrido algo reseñable desde el último seguimiento?} \checkedbox\ Sí · \checkbox\ No

Tres semanas desde el seguimiento del S9.
Del 26-10 al 11-11 no hubo sesiones, por exámenes y otros compromisos, y el 12-11 el método se simplificó: el \eng{sprint} pasó a ser una tanda de trabajo, el seguimiento semanal se sustituyó por este control mensual y el plan se reemitió como versión 8.
La auditoría de coherencia del 12-11 encontró los restos del método anterior en la memoria y en las instrucciones, dos datos privados y un perfil de permisos demasiado amplio; todo quedó corregido el 15-11.

\Needspace*{6\baselineskip}
\textbf{Nuevos riesgos identificados} \updatesymbol

\begin{tabularx}{\textwidth}{|C{1.4cm}|Y|C{0.8cm}|C{0.8cm}|C{1cm}|C{2.4cm}|}
\hline
\multicolumn{1}{|C{1.4cm}|}{\textbf{Id}} &
\multicolumn{1}{>{\centering\arraybackslash}X|}{\textbf{Descripción}} &
\multicolumn{1}{C{0.8cm}|}{\textbf{P}} &
\multicolumn{1}{C{0.8cm}|}{\textbf{I}} &
\multicolumn{1}{C{1cm}|}{\textbf{Exp.}} &
\multicolumn{1}{C{2.4cm}|}{\textbf{Severidad}} \\
\hline
— & Ninguno. & — & — & — & — \\
\hline
\end{tabularx}

\textbf{¿Afecta algo a los objetivos o al alcance (cap.~3 de la memoria)?} \checkbox\ Sí · \checkedbox\ No

\textbf{Justificación:} el retraso de tres semanas lo absorbe la holgura del plan; el alcance no cambia.

\textbf{Otros eventos relevantes:} el perfil de permisos de la auditoría de coherencia se endureció (sin órdenes que permitan leer los archivos de secretos ni escribir en GitHub); gasto de noviembre en la suscripción de IA, 9~€, con el plan Pro (cap.~14).

\end{document}
